\chapter{Introducción}
\label{cap:introduccion}

\section{Motivación}

Los modelos de lenguaje dejaron de ser herramientas que sólo completan líneas
de código. Los agentes actuales pueden recorrer un repositorio, ejecutar
comandos, interpretar errores y modificar varios archivos en ciclos sucesivos
de observación y acción \cite{chen2021codex,yao2023react,yang2024sweagent}. Esta
evolución amplió el tipo de problemas que pueden abordar, pero también volvió
visible una diferencia central: resolver un parche local no es equivalente a
coordinar un cambio que atraviesa dominio, persistencia, interfaces y pruebas.

Un agente único puede perder información relevante cuando debe sostener a la
vez demasiados archivos, dependencias y decisiones. Dividir el trabajo entre
varios agentes parece una solución natural, pero introduce un problema nuevo:
cada frontera creada debe coordinarse y luego integrarse. Dos agentes pueden
implementar firmas incompatibles, editar el mismo archivo o aprobar pruebas que
no demuestran el comportamiento solicitado. Por lo tanto, el desafío no es
simplemente utilizar más agentes, sino gobernar su colaboración como un proceso
de ingeniería de software.

Esta tesis aborda ese desafío mediante \sistema, un plano de control local que
coordina agentes de código sin reemplazarlos. Los modelos conservan aquello en
lo que son valiosos —interpretar objetivos y proponer implementaciones—,
mientras que el sistema determina qué puede ejecutarse, qué debe esperar, qué
rutas puede modificar cada intento y qué evidencia se exige antes de entregar
un resultado.

\section{El problema: la paradoja de la granularidad}

La modularidad busca dividir un sistema alrededor de responsabilidades que
puedan cambiar de forma relativamente independiente \cite{parnas1972criteria}.
En un sistema multi-agente, esta idea adquiere una consecuencia operacional:
la forma de dividir el objetivo determina cuántos agentes se ejecutan, cuánto
contexto recibe cada uno y cuántas integraciones deben resolverse.

Existen dos extremos igualmente problemáticos:

\begin{itemize}
  \item \textbf{Una tarea demasiado grande.} Un solo agente debe comprender
  muchas decisiones y archivos. Aumentan la duración del intento, la pérdida de
  contexto y las modificaciones fuera de alcance.
  \item \textbf{Una tarea excesivamente fragmentada.} Cada fragmento requiere
  contexto, contrato, proceso y validación propios. Crecen las costuras entre
  partes y, con ellas, el riesgo de incompatibilidades y el costo de integración
  \cite{brooks1995mythical}.
\end{itemize}

Se denomina \textbf{paradoja de la granularidad} a la necesidad de encontrar
una división suficientemente pequeña para ser ejecutable, pero suficientemente
cohesiva para que la coordinación no domine el trabajo útil. No existe un
tamaño universal: la decisión depende de las fronteras del repositorio, de los
artefactos compartidos y de la posibilidad de verificar cada parte.

\section{Propuesta de solución}

ManyHands trata un objetivo de software como un \emph{run}: una transformación
completa desde un repositorio base hasta un commit integrado y entregado. El
run se organiza mediante cinco estrategias complementarias:

\begin{enumerate}
  \item \textbf{Grafo híbrido.} La jerarquía expresa quién integra a quién; las
  relaciones tipadas expresan dependencias de artefactos, acuerdos de interfaz
  y conflictos de escritura.
  \item \textbf{Contratos versionados.} Cada unidad conoce su objetivo, alcance,
  artefactos, interfaces y obligaciones de validación antes de ejecutar.
  \item \textbf{Granularidad explicable.} Una división se conserva sólo si
  aporta capacidad, paralelismo o aislamiento de pruebas.
  \item \textbf{Aislamiento y recuperación.} Los intentos se ejecutan en
  worktrees separados, bajo supervisión de procesos y con estado durable.
  \item \textbf{Evidencia antes que declaración.} La integración y la entrega
  dependen de pruebas y recibos asociados al commit exacto, no del relato del
  modelo.
\end{enumerate}

Estas estrategias forman una cadena de custodia: cada etapa recibe una entrada
identificada, produce un artefacto verificable y deja una evidencia durable para
la etapa siguiente.

\section{Preguntas y objetivos de investigación}

El trabajo se organiza alrededor de dos preguntas:

\begin{description}[leftmargin=1.8cm,style=nextline]
  \item[\textbf{PI-1}] ¿Es viable coordinar agentes de código mediante un grafo
  y contratos versionados, aislando su ejecución e integrando sus resultados de
  manera que cada entrega quede vinculada a evidencia reproducible sobre un
  commit exacto?
  \item[\textbf{PI-2}] ¿Puede una política determinista y explicable decidir
  cuándo dividir una tarea o conservarla como una unidad cohesiva a partir de
  propiedades observables de capacidad, paralelismo y verificación?
\end{description}

La hipótesis principal sostiene que una arquitectura basada en contratos,
relaciones tipadas, aislamiento físico e integración por evidencia permite
coordinar agentes heterogéneos sin depender de correcciones imperativas del
grafo. La hipótesis complementaria sostiene que una política categórica puede
evitar la microfragmentación sin recurrir a puntajes difíciles de interpretar.

El \textbf{objetivo general} es diseñar, implementar y evaluar una plataforma
local de orquestación multi-agente que transforme objetivos de software en
resultados integrados, auditables y reproducibles. De él se derivan los
siguientes objetivos específicos:

\begin{enumerate}
  \item modelar el trabajo mediante grafos jerárquicos y relaciones tipadas;
  \item formalizar contratos de objetivo, alcance, interfaz, artefacto y
  validación;
  \item diseñar una política categórica de granularidad y un algoritmo de
  colapso que preserve criterios;
  \item implementar ejecución aislada, supervisión de procesos y recuperación
  durable;
  \item integrar resultados de forma ascendente y construir evidencia sobre el
  candidato exacto;
  \item publicar entregas transaccionales y conservar un historial auditable;
  \item evaluar el ciclo completo mediante regresiones y un caso experimental
  representativo.
\end{enumerate}

\section{Alcance}

ManyHands se diseñó como una herramienta local para un desarrollador. Opera
sobre repositorios Git y utiliza ejecutores externos de modelos mediante CLI.
El trabajo prioriza corrección, contención y trazabilidad por encima de la
velocidad máxima. Quedan fuera de alcance la colaboración multiusuario, el
entrenamiento de modelos, la optimización de costos entre proveedores y el
aislamiento mediante máquinas virtuales completas.

La evaluación estudia la viabilidad de la arquitectura y su capacidad para
producir una descomposición útil. No pretende establecer una comparación
estadística universal entre modelos ni generalizar rendimiento a todos los
lenguajes y tamaños de repositorio.

\section{Organización del documento}

Los capítulos~\ref{cap:preliminares} y~\ref{cap:estado} introducen los conceptos
y antecedentes. El capítulo~\ref{cap:metodologia} presenta la política de
granularidad. Los capítulos~\ref{cap:arquitectura} y~\ref{cap:implementacion}
explican el diseño y su materialización. El capítulo~\ref{cap:evaluacion}
describe el experimento \emph{Viaje en Familia}; el
capítulo~\ref{cap:discusion} analiza sus alcances y el
capítulo~\ref{cap:conclusiones} sintetiza los aportes y líneas futuras.
